\documentclass[conference]{IEEEtran}
\usepackage{times}
\usepackage[numbers]{natbib}
\usepackage[bookmarks=true]{hyperref}

\pdfinfo{
   /Author (Zhijie Chen and Shiyi Xiao)
   /Title  (Dynamic-Obstacle Quadruped Parkour with Robust Depth Perception)
   /CreationDate (D:20260420120000)
   /Subject (Intelligent Robot Project Proposal)
   /Keywords (quadruped parkour; reinforcement learning; dynamic obstacles)
}

\begin{document}

\title{Dynamic-Obstacle Quadruped Parkour with\\Robust Depth Perception}

\author{
   \IEEEauthorblockN{Zhijie Chen\IEEEauthorrefmark{1} \and Shiyi Xiao\IEEEauthorrefmark{1}}
   \IEEEauthorblockA{\IEEEauthorrefmark{1}Shanghai Jiao Tong University\\
      \{zhijie.chen, sjtuxsy-finance-pc\}@sjtu.edu.cn\\
      Equal contribution}
}

\maketitle

\begin{abstract}
   We propose a research project on quadruped parkour that builds on \emph{Extreme Parkour with Legged Robots}~\citep{cheng2023parkour}. We will first reproduce its simulation training pipeline and validate performance on standard static obstacle tracks, as required by the course project. We will then extend the benchmark from static to dynamic obstacles, including moving hurdles, shifting gaps, time-varying ramps, and changing step heights. Our main technical idea is to augment the original two-phase pipeline with latent variables that encode dynamic obstacle state, so that a depth-based deployment policy can better infer time-varying geometry for action timing and foot placement. As a secondary, higher-risk direction, we will study whether short temporal fusion and sensing degradations during training improve robustness to low-frequency, jittery depth observations. We expect the baseline policy to degrade on dynamic obstacles, and we aim to show that the modified pipeline improves performance on at least some visually recoverable dynamic settings.
\end{abstract}

\section{Motivation and Problem Statement}
This project is a research extension of the quadruped parkour topic in Section 2.4 of the course idea pool. The baseline paper~\citep{cheng2023parkour} shows that a Unitree A1 robot can traverse hurdles, gaps, tilted ramps, and high steps with a policy trained from privileged simulation inputs and distilled to a deployable depth-based policy. This provides a strong starting point for reproduction, but its benchmark is still built around static terrain geometry. We want to study whether this learning framework can handle a harder and more realistic setting in which obstacle geometry changes over time.

The project therefore has two parts. The first is required baseline work: reading the paper and codebase, reproducing the training setup, and validating the policy on standard static tracks. The second is our proposed extension. The main extension is dynamic-obstacle parkour, where successful traversal depends on both obstacle geometry and obstacle state at execution time. A secondary extension is robustness under degraded depth sensing, motivated by the baseline paper's own discussion of low-frequency, jittery, and artifact-prone perception.

\section{Literature Review}
Recent legged locomotion work has shown that policies trained with privileged information and deployed from onboard sensing can handle challenging terrain using adaptation and vision~\citep{kumar2021rma,agarwal2022legged,miki2022learning}. The parkour baseline of \citet{cheng2023parkour} combines this line of work with a two-phase training strategy: Phase 1 uses privileged exteroception and waypoint-derived heading, while Phase 2 distills the policy to depth input and predicted heading. This makes it a suitable baseline for our project because it already provides a strong parkour pipeline, a clear codebase, and an explicit limitation: the paper contrasts learning-based control with methods that assume fixed obstacles and directly asks what happens when obstacles move.

Our proposal follows that opening. Instead of designing a new parkour task from scratch, we retain the original pipeline and ask whether it can be extended to time-varying terrain. We also keep a smaller robustness branch because the same paper emphasizes that depth sensing is noisy, low-frequency, and jittery, suggesting a possible bottleneck at deployment.

\section{Technical Plan}
We will use the official \texttt{extreme-parkour} implementation as the baseline codebase. First, we will reproduce the original simulation policy and validate performance on standard static obstacles. This covers the reproduction and validation requirements of the course project.

Our main extension is dynamic-obstacle parkour. We will convert originally static terrain elements into dynamic ones: hurdles that move back and forth, gaps whose effective takeoff or landing geometry shifts over time, ramps whose slope changes, and steps whose height changes. Compared with the static setting, this extension requires the policy to estimate not only obstacle geometry but also obstacle state at the moment of execution. We will compare three settings: static obstacles with the original pipeline, dynamic obstacles with the original pipeline, and dynamic obstacles with an improved pipeline. The improved pipeline will augment the latent environment information with dynamic obstacle state, such as current hurdle position, gap configuration, step height, or ramp angle. In Phase 2, the depth-based policy will be trained to recover these quantities explicitly or implicitly from observations. We will decide empirically whether ROA-style latent estimation, teacher-student distillation, or a hybrid version works better for this extension, based on both performance and training stability.

Our secondary branch is robust depth perception under degraded sensing. Here we will test a lightweight extension consisting of short-horizon temporal fusion, frame-drop augmentation, and latency/jitter augmentation during training. This branch is intentionally smaller and higher risk than the dynamic-obstacle work; its purpose is to test whether modest perception-side changes improve robustness rather than to redesign the entire control pipeline.

\section{Expected Outcomes}
For the main branch, we expect the original pipeline to perform worse once static obstacles are replaced by dynamic ones. We then expect the improved pipeline to recover part of this loss on at least some dynamic obstacle types, especially when the obstacle state is visually recoverable from egocentric depth. We will evaluate success rate, takeoff position error, landing position error, and representative failure modes such as late takeoff, obstacle collision, and unstable landing.

For the secondary branch, the expected outcome is more modest. We will measure whether the modified perception stack is less sensitive to lower frame rates, frame drops, latency perturbation, and depth corruption. Even if the improvement is small, the experiments can still clarify how fragile the baseline depth pipeline is under degraded sensing.

\section{Potential Risks and Backup Plan}
The main risk is that some dynamic obstacle states may be weakly observable from the onboard depth view, which would limit the value of the proposed latent-variable extension. If this happens, we will narrow the scope to the most feasible dynamic obstacle types, such as those whose motion is more directly visible, and present a careful analysis of where the method helps and where it does not.

The secondary branch is higher risk. It may produce only small gains, or progress on it may be limited by time after completing the main branch. Our backup plan is therefore to prioritize the baseline reproduction and the dynamic-obstacle experiments, while treating robust perception as a smaller ablation-oriented branch. Even partial results from this branch remain useful for the proposal because they inform limitations, stress testing, and failure analysis.

\section{Timeline}
The semester is currently in Week 8, so the remaining work is scheduled by week:
\begin{itemize}
   \setlength{\itemsep}{0.15em}
   \item Weeks 8--9: read the baseline paper and codebase, set up the reproduction workflow, and validate the baseline on static obstacle tracks.
   \item Weeks 10--12: implement dynamic terrains and run the main training and evaluation for the dynamic-obstacle extension.
   \item Weeks 13--14: work on the robustness branch if progress permits, with the expectation that this part may take only one week if results are weak or time is tight.
   \item Weeks 15--16: organize experimental results and write the final report; if the robustness branch finishes early or is deprioritized, one additional week will be shifted here.
\end{itemize}

\bibliographystyle{plainnat}
\bibliography{references}

\end{document}
